% TEMPLATE-MILC-v3.tex - plantilla maestra UNICA de las guias MILC v3.
%
% La usan las tres familias de guias, cada una con su driver:
%   · Grados 6-11  -> scripts/build-guias-g11.py
%   · Bebras       -> scripts/build-guias-bebras.py
%   · Semillero    -> scripts/build-guias-semillero.py
%
% Antes habia tres copias casi identicas (~400 lineas iguales, 6 distintas):
% cada mejora habia que aplicarla tres veces, y en la practica no se hacia
% (el motor de recursos vivia solo en dos de ellas). Lo que varia entre
% familias esta parametrizado en 6 placeholders que cada driver llena
% (se nombran aqui SIN su sintaxis real: el builder reemplaza por texto plano,
% tambien dentro de los comentarios, y un valor multilinea rompe el preambulo):
%
%   HEADER_IZQ        encabezado izquierdo de pagina
%   HEADER_DER        encabezado derecho de pagina
%   PORTADA_ETIQUETA  rotulo sobre el numero grande (GRADO/MOMENTO/MODULO)
%   PORTADA_SERIE     linea de serie de la portada
%   PORTADA_ID        identificador de la guia en la portada
%   PORTADA_CIRCULO   el circulito de grado (°), o vacio si no aplica
%
% El resto de claves las documenta content/guias/_SCHEMA.md. El script valida
% que no queden placeholders sin reemplazar antes de escribir el archivo final.

\documentclass[11pt,letterpaper]{article}
\usepackage[margin=2.0cm]{geometry}
\usepackage[table]{xcolor}
\usepackage{fontspec}
\usepackage[spanish]{babel}
\usepackage{tikz}
% Librerías TikZ para los diagramas de las guías (bloque `recursos:`):
%  · babel  → babel en español vuelve activos caracteres como «>», y TikZ los usa
%             en sus opciones (>=stealth, ->). Sin esta librería, cualquier
%             diagrama con flechas revienta con «\language@active@arg> has an
%             extra }». Debe ir ANTES de que se dibuje nada.
%  · positioning        → sintaxis «below left=2pt and 3pt».
%  · decorations.pathreplacing → llaves ({brace}) para señalar rangos.
\usetikzlibrary{babel,positioning,decorations.pathreplacing}
\usepackage{graphicx}
% Circuitos: el trabajo de electrónica se hace hoy con micro:bit y sus sensores
% integrados (MakeCode, sin cableado), así que no se necesitan esquemáticos.
% Si algún día entran montajes con protoboard, basta descomentar esta línea:
% los diagramas .tex/.tikz declarados en `recursos:` ya se incrustan solos.
% \usepackage{circuitikz}
\usepackage[most]{tcolorbox}
\usepackage{tabularx}
\usepackage{array}
\usepackage{fancyhdr}
\usepackage{titlesec}
\usepackage[document]{ragged2e}  % bandera a la derecha en todo el cuerpo
\usepackage{enumitem}
\usepackage{microtype}

% Helvetica Neue no trae la flecha U+2192, asi que cada «→» escrito literalmente
% en el YAML se compilaba a NADA, en silencio (462 flechas en 61 guias: rutas del
% tipo «paleta → tipografia → lockup» salian rotas). Mapeamos el caracter a su
% version matematica, que si tiene glifo. Asi el autor puede seguir escribiendo
% «→» comodamente en el YAML.
\usepackage{newunicodechar}
\newunicodechar{→}{\ensuremath{\rightarrow}}
% Mismo problema con los simbolos de marca y vinetas: el «✓» de «una palomita ✓»
% salia como cuadrito vacio en el PDF de todas las guias que lo usaban.
\usepackage{pifont}
\newunicodechar{✓}{\ding{51}}
\newunicodechar{✗}{\ding{55}}
\newunicodechar{✔}{\ding{52}}
\newunicodechar{•}{\textbullet}
\newunicodechar{≥}{\ensuremath{\geq}}
\newunicodechar{≤}{\ensuremath{\leq}}

\setmainfont{Helvetica Neue}
\setsansfont{Helvetica Neue}
\setlength{\parindent}{0pt}
\setlength{\parskip}{6pt}
% Sin particion de palabras: en 6.º-7.º una palabra cortada por guion al final
% de linea («Comuni-cacion») frena la lectura mucho mas que una linea corta.
\hyphenpenalty=10000
\exhyphenpenalty=10000
\pretolerance=2000
\tolerance=3000
\setlength{\emergencystretch}{3em}
\linespread{1.10}
\renewcommand{\arraystretch}{1.25}
\setlist{leftmargin=5mm,itemsep=1mm,topsep=1mm}
\newcolumntype{Y}{>{\RaggedRight\arraybackslash}X}

\definecolor{milcMagenta}{HTML}{E6007E}
\definecolor{milcVino}{HTML}{5A0038}
\definecolor{milcTurquesa}{HTML}{0093A5}
\definecolor{milcVerde}{HTML}{58A923}
\definecolor{milcMostaza}{HTML}{E5B400}
\definecolor{milcOcre}{HTML}{B86600}
\definecolor{milcNegro}{HTML}{241816}
\definecolor{milcGris}{HTML}{F7F7F4}
\definecolor{milcLinea}{HTML}{E8E2DD}
\definecolor{milcGradePrimary}{HTML}{84CC16}
\definecolor{milcGradeDark}{HTML}{4D7C0F}
\definecolor{milcGradeSoft}{HTML}{ECFCCB}
\definecolor{milcGradeAccent}{HTML}{CFE7FF}
\definecolor{milcGradeText}{HTML}{FFFFFF}
\color{milcNegro}

\pagestyle{fancy}
\fancyhf{}
\lhead{\small Tecnología e Informática · Bebras}
\rhead{\small Momento 8 · MILC v3}
\cfoot{\small\thepage}
\renewcommand{\headrulewidth}{0pt}

\newtcolorbox{titlebox}[1]{
  enhanced,colback=#1,colframe=#1,arc=7mm,boxrule=0pt,
  left=7mm,right=7mm,top=3mm,bottom=3mm,
  before skip=8pt,after skip=8pt,
  fontupper=\sffamily\bfseries\Large\color{white}
}

\newtcolorbox{softbox}[3]{
  enhanced,colback=#2,colframe=#1,borderline west={4pt}{0pt}{#1},
  arc=2mm,boxrule=.4pt,left=4mm,right=4mm,top=3mm,bottom=3mm,
  before skip=6pt,after skip=8pt,title={#3},coltitle=white,
  colbacktitle=#1,fonttitle=\sffamily\bfseries,fontupper=\RaggedRight,
  attach boxed title to top left={xshift=3mm,yshift=-2mm},
  boxed title style={arc=2mm, boxrule=0pt}
}

\newtcolorbox{drawbox}[2]{
  enhanced,colback=white,colframe=#1,arc=4mm,boxrule=1pt,
  left=5mm,right=5mm,top=4mm,bottom=4mm,before skip=8pt,after skip=8pt,
  title={#2},coltitle=white,colbacktitle=#1,fonttitle=\sffamily\bfseries,
  fontupper=\RaggedRight,
  attach boxed title to top left={xshift=4mm,yshift=-2mm},
  boxed title style={arc=3mm, boxrule=0pt}
}

\newtcolorbox{infoband}[1]{
  enhanced,colback=#1!8,colframe=#1!50,arc=2mm,boxrule=.5pt,
  left=4mm,right=4mm,top=2mm,bottom=2mm,before skip=4pt,after skip=4pt,
  fontupper=\small\RaggedRight
}

% Puente narrativo entre secciones (contrato editorial Conéctate).
% Da continuidad: "Acabas de X. Ahora vas a Y porque Z."
% Visualmente: banda izquierda en color del grado, texto en itálica pequeña.
\newtcolorbox{puentebox}{
  enhanced,colback=milcGradeSoft,colframe=milcGradeDark,
  borderline west={3pt}{0pt}{milcGradeDark},
  arc=0mm,boxrule=0pt,left=5mm,right=4mm,top=2.5mm,bottom=2.5mm,
  before skip=4pt,after skip=8pt,
  fontupper=\itshape\small\color{milcNegro}\RaggedRight
}
% La flecha va en modo matematico a proposito: Helvetica Neue NO tiene el glifo
% U+2192, asi que \textrightarrow se compilaba a NADA («Missing character: There
% is no → in font Helvetica Neue») y el puente salia sin flecha. En modo
% matematico la flecha viene de la fuente matematica y si se dibuja.
\newcommand{\puente}[1]{%
  \begin{puentebox}%
    \textcolor{milcGradeDark}{$\rightarrow$}\hspace{2mm}#1%
  \end{puentebox}%
}

\newcommand{\linead}[1]{\noindent\color{milcLinea}\rule{#1}{0.5pt}\color{milcNegro}}
\newcommand{\respuesta}[1]{\par\vspace{2mm}\foreach \n in {1,...,#1}{\noindent\color{milcLinea}\rule{\linewidth}{0.5pt}\par\vspace{5mm}}\color{milcNegro}}
\newcommand{\checkbox}{\textcolor{milcGradeDark}{\fbox{\phantom{x}}}\hspace{5pt}}

\newcommand{\phasecard}[4]{
\begin{tcolorbox}[enhanced,colback=#3,colframe=#2,arc=3mm,boxrule=.5pt,height=3.05cm,valign=center,halign=center,left=1.5mm,right=1.5mm,top=2mm,bottom=2mm]
{\sffamily\bfseries\fontsize{22}{24}\selectfont\textcolor{#2}{#1}}\par
{\sffamily\bfseries\small #4}
\end{tcolorbox}
}

% ── Piezas del rediseno 2026 (legibilidad 6.º-11.º) ───────────────────
% Banda de estacion: reemplaza el titlebox macizo. Titulo a la izquierda,
% dato practico (tiempo/modalidad) a la derecha, en una sola linea.
\newcommand{\milcestacion}[2]{%
\begin{tcolorbox}[enhanced,colback=#1,colframe=#1,arc=2mm,boxrule=0pt,
  left=5mm,right=5mm,top=2.6mm,bottom=2.6mm,before skip=11pt,after skip=7pt]
{\sffamily\bfseries\fontsize{15}{18}\selectfont\color{white}#2}
\end{tcolorbox}}

% Tarjeta de paso numerada: sustituye las tablas «Pilar / Aplicacion», que
% eran formato de consulta docente, no de estudiante. El numero va en un
% circulo sobre el borde para que la secuencia se lea de un vistazo.
\newcommand{\milcpaso}[3]{%
\begin{tcolorbox}[enhanced,colback=white,colframe=milcLinea,arc=1.5mm,boxrule=.9pt,
  left=14mm,right=4mm,top=3mm,bottom=3mm,before skip=4pt,after skip=4pt,
  fontupper=\RaggedRight,
  overlay={\node[anchor=north west,circle,fill=#2,text=white,inner sep=0pt,
    minimum size=7.4mm,font=\sffamily\bfseries\small]
    at ([xshift=3.6mm,yshift=-3.2mm]frame.north west) {#1};}]
#3
\end{tcolorbox}}

% Casilla marcable de tamano util con lapiz (la anterior era \fbox{\phantom{x}},
% de alto variable segun el texto que la seguia).
\newcommand{\milcmarca}{\framebox[3.8mm]{\rule{0pt}{3.4mm}}}

% Fila marcable de lista suelta (checks de la estacion 1).
\newcommand{\milccheckrow}[1]{%
\par\vspace{1.8mm}\noindent
\makebox[6.4mm][l]{\raisebox{-.55mm}{\textcolor{milcNegro}{\milcmarca}}}%
\parbox[t]{\dimexpr\linewidth-6.4mm\relax}{\RaggedRight #1}\par}

% Celda marcable dentro de la rubrica (tabla de autoevaluacion). Misma
% sangria francesa, pero pensada para vivir dentro de una columna X.
\newcommand{\milcopcion}[1]{%
\makebox[5.2mm][l]{\raisebox{-.55mm}{\textcolor{milcNegro}{\milcmarca}}}%
\parbox[t]{\dimexpr\linewidth-5.2mm\relax}{\RaggedRight #1}}

% ── Recursos visuales de la guía ──────────────────────────────────────
% Declarados en el bloque `recursos:` del YAML y emitidos por el builder.
% \guiaFigura   → imagen rasterizada o PDF (foto, carta celeste, montaje).
% \guiaDiagrama → fuente TikZ/circuitikz (.tex) incrustada como vector:
%                 es la vía para circuitos y esquemáticos editables.
% #1 = ruta relativa al .tex · #2 = pie (puede ir vacío)
\newcommand{\guiaFigura}[2]{%
\begin{center}
\includegraphics[width=0.88\linewidth,height=0.38\textheight,keepaspectratio]{#1}\\[1.5mm]
{\footnotesize\itshape\textcolor{milcNegro!75}{#2}}
\end{center}
}

\newcommand{\guiaDiagrama}[2]{%
\begin{center}
\input{#1}\\[1.5mm]
{\footnotesize\itshape\textcolor{milcNegro!75}{#2}}
\end{center}
}

\begin{document}
\thispagestyle{empty}

% ── Portada ───────────────────────────────────────────────────────────
% Antes: un rectangulo de color a pagina completa. Se veia bien en pantalla,
% pero en la fotocopiadora del colegio se comia el toner y salia gris sucio.
% Ahora la banda ocupa el tercio superior y el resto va en blanco.
\begin{tikzpicture}[remember picture,overlay]
  \path[fill=milcGradePrimary]
    (current page.north west) rectangle ([yshift=-10.6cm]current page.north east);

  \node[anchor=north west,text=milcGradeText] at ([xshift=2.0cm,yshift=-1.55cm]current page.north west)
    {\sffamily\bfseries\fontsize{12}{14}\selectfont MOMENTO};
  \node[anchor=north west,text=milcGradeText,opacity=.85] at ([xshift=2.0cm,yshift=-2.35cm]current page.north west)
    {\sffamily\bfseries\fontsize{8.5}{10}\selectfont BEBRAS · PENSAR COMO UNA MÁQUINA};
  \node[anchor=north east,text=milcGradeText,opacity=.85,align=right] at ([xshift=-2.0cm,yshift=-1.55cm]current page.north east)
    {\sffamily\bfseries\fontsize{9}{12}\selectfont I.E. Sor María Juliana\\Cartago, Valle del Cauca};

  \node[anchor=west,text=milcGradeText] (gradonum) at ([xshift=1.85cm,yshift=-7.1cm]current page.north west)
    {\sffamily\bfseries\fontsize{112}{112}\selectfont 8};
  % El circulito de grado (°) solo aplica a las guias de grado («11°»); en
  % Bebras y Semillero el numero grande es un momento/modulo, no un grado.
  % Se posiciona relativo al nodo `gradonum`, no en coordenadas absolutas.
  

  \node[anchor=west,text=milcGradeText,text width=11.4cm,align=left] at ([xshift=1.5cm,yshift=0.5cm]gradonum.east)
    {
     {\sffamily\bfseries\fontsize{9}{11}\selectfont\MakeUppercase{Curso «Pensar como una máquina» · Pensamiento computacional}}\\[3mm]
     {\sffamily\bfseries\fontsize{21}{25}\selectfont\setlength{\baselineskip}{27pt}Modo concurso --- estrategia de ataque,\\[4pt]tiempo y puntuación}\\[3.5mm]
     {\sffamily\bfseries\fontsize{15}{18}\selectfont Bebras · Momento 8}
    };
\end{tikzpicture}

\vspace*{9.4cm}
\vfill

{\sffamily\bfseries\fontsize{8}{10}\selectfont\textcolor{milcGradeDark}{LO QUE TE LLEVAS HOY}}

\begin{tcolorbox}[enhanced,colback=white,colframe=milcNegro,arc=2mm,boxrule=1.6pt,
  left=5mm,right=5mm,top=4mm,bottom=4mm,before skip=3pt,after skip=12pt,
  fontupper=\RaggedRight\fontsize{11}{15.5}\selectfont]
Tu tabla de decisión estratégica (escenario con 3 tareas justificado con valor esperado) más tu plan de ataque personal para los 45 minutos.
\end{tcolorbox}

\begin{tcolorbox}[enhanced,colback=milcGris,colframe=milcGris,arc=2mm,boxrule=0pt,
  left=5mm,right=5mm,top=3.5mm,bottom=3.5mm,after skip=12pt]
{\sffamily\bfseries\fontsize{8}{10}\selectfont\textcolor{milcNegro!55}{ESTA GUÍA ES DE}}\\[3mm]
\begin{tabularx}{\linewidth}{X p{3.2cm} p{3.2cm}}
\linead{\linewidth} & \linead{3.0cm} & \linead{3.0cm}\\[-1mm]
{\fontsize{8.5}{10}\selectfont\textcolor{milcNegro!55}{Nombre completo}} &
{\fontsize{8.5}{10}\selectfont\textcolor{milcNegro!55}{Grupo}} &
{\fontsize{8.5}{10}\selectfont\textcolor{milcNegro!55}{Fecha}}\\
\end{tabularx}
\end{tcolorbox}

\vfill

\noindent\textcolor{milcLinea}{\rule{\linewidth}{1pt}}\\[2mm]
{\sffamily\bfseries\fontsize{9.5}{12}\selectfont Autor: Dr. Álvaro Cárdenas Orozco}\\[1mm]
{\sffamily\fontsize{8.5}{11}\selectfont\textcolor{milcNegro!55}{Metodología MILC · apertura ancestral · pensamiento computacional · triángulo Dussel–estoicismo–Floridi}}

\newpage

\section*{Apertura: Modo concurso: estrategia de ataque, tiempo y puntuación}

\begin{softbox}{milcGradeDark}{milcGradeSoft}{Saber ancestral}
En los barrios de Cartago y en las calles del Valle del Cauca, las \textbf{canicas y el trompo} no eran solo diversión: eran escuela de estrategia. El buen jugador de canicas no tira al azar. Sabe cuál bolita arriesgar y cuál guardar. Estudia el terreno, lee la posición del contrario y decide con calma: «si tiro aquí, ¿qué gano?, ¿qué arriesgo?». Con el trompo pasa lo mismo: antes de lanzarlo, el experto lee el piso, elige el ángulo y sabe que un mal lanzamiento puede costar la peona más valiosa. Esa inteligencia táctica --- arriesgar con cabeza, no por impulso --- no estaba en ningún libro. La enseñaba la calle, la transmitía el vecino mayor, la perfeccionaba el juego mismo. \emph{Origen: juegos tradicionales de calle colombianos, cultura popular del Valle del Cauca.}
\end{softbox}

\begin{softbox}{milcTurquesa}{milcGris}{Saber contemporáneo}
\textbf{Bebras} tiene tres niveles de dificultad --- A (fácil), B (media) y C (difícil) --- y cada uno tiene su propia \textbf{puntuación}: A da $+6$ si aciertas, $0$ si fallas o dejas en blanco; B da $+9$ si aciertas, $-2$ si fallas, $0$ si dejas en blanco; C da $+12$ si aciertas, $-4$ si fallas, $0$ si dejas en blanco. Esa diferencia lo cambia todo: en A nunca pierde quien responde aunque adivine, mientras que en B y C adivinar al azar sin descartar opciones te perjudica en promedio. Saber la respuesta no es suficiente; hay que saber \emph{cuándo responder}, cuándo arriesgar y cuándo es mejor dejar en blanco. La \textbf{penalización} convierte el concurso en un problema de toma de decisiones, no solo de conocimiento.
\end{softbox}

\begin{softbox}{milcMostaza}{milcGris}{Pregunta-puente}
\textit{En los juegos de calle, ganar no dependía solo de habilidad: dependía de saber cuándo arriesgar y cuándo parar. ¿Y si en Bebras también ganaras no solo por saber la respuesta, sino por decidir bien cuándo responder, cuándo arriesgar y cuándo pasar?}
\end{softbox}

\begin{softbox}{milcVerde}{milcGris}{Saber y saber hacer}
Al terminar podrás: (1) \textbf{identificar} la tabla de puntuación oficial de Bebras (A: $+6/0$, B: $+9/-2$, C: $+12/-4$) y explicar qué cambia con la penalización; (2) \textbf{analizar} un escenario de concurso y calcular el valor esperado de cada decisión (responder, adivinar, dejar en blanco); (3) \textbf{evaluar} la mejor jugada en una situación con tiempo limitado y tres tareas pendientes, justificando con números; (4) \textbf{aplicar} esa lógica en tu propio plan de ataque para los 45 minutos del concurso.
\end{softbox}



\small
\begin{tabularx}{\textwidth}{>{\bfseries}p{3.0cm}Y}
\rowcolor{milcGradePrimary!12}Autor & Dr. Álvaro Cárdenas Orozco\\
DBA & Desarrolla el pensamiento computacional --- descomposición, patrones, abstracción y algoritmos --- para resolver problemas (transversal 6°--11°, alineado a los lineamientos MEN de Tecnología e Informática)\\
Referentes & Valentina Dagienė (Bebras) · Pensamiento computacional (Wing) · Dussel · Estoicismo · Floridi · Saberes del Valle y el Pacífico\\
Producto final & Tu tabla de decisión estratégica (escenario con 3 tareas justificado con valor esperado) más tu plan de ataque personal para los 45 minutos.\\
\end{tabularx}
\normalsize

\puente{Ya sabes que las canicas enseñan a medir el riesgo antes de lanzar. Ahora mira el mapa de la sesión: vas a aprender las reglas de puntuación, a calcular cuándo conviene arriesgar y a diseñar tu plan para los 45 minutos.}

\milcestacion{milcGradeDark}{Ruta MILC: así vamos a aprender}
\begin{tabularx}{\textwidth}{>{\centering\arraybackslash}X>{\centering\arraybackslash}X>{\centering\arraybackslash}X>{\centering\arraybackslash}X}
\phasecard{1}{milcVerde}{milcGris}{Escucha\\[-1mm]\small Observo, escucho y pregunto.} &
\phasecard{2}{milcTurquesa}{milcGris}{Entender\\[-1mm]\small Sistematizo y ordeno.} &
\phasecard{3}{milcMagenta}{milcGris}{Hacer\\[-1mm]\small Construyo y aplico.} &
\phasecard{4}{milcVino}{milcGris}{Cierre\\[-1mm]\small Evalúo y reflexiono.}
\end{tabularx}


\puente{Antes de ver los números, conviene que recuerdes cómo tomabas decisiones bajo presión en un juego. Empieza observando una situación donde el riesgo y la recompensa se pesan antes de actuar.}

\milcestacion{milcVerde}{Estación 1 de 3 · Escucha}

\begin{softbox}{milcVerde}{milcGris}{Escena para escuchar}
\textbf{Recuerda un juego de canicas o trompo} --- o cualquier juego de calle en el que hayas jugado con alguien de tu barrio. Piensa en un momento donde tuviste que \emph{decidir si arriesgabas} algo valioso o lo guardabas. ¿Qué información tenías antes de decidir? ¿Qué podías perder? ¿Qué podías ganar? ¿Cómo calculaste si valía la pena? Quizás no lo hiciste con fórmulas, pero sí con una lógica: «si lanzo aquí, puedo ganar tres y perder una; eso me conviene». Esa misma lógica rige la estrategia de Bebras.
\end{softbox}

{\sffamily\bfseries\fontsize{8}{10}\selectfont\textcolor{milcNegro!55}{MARCA CUANDO LO HAYAS HECHO}}
\milccheckrow{Nombraste al menos una situación concreta donde decidiste arriesgar o no en un juego.}
\milccheckrow{Identificaste qué podías ganar y qué podías perder en esa situación.}
\milccheckrow{Señalaste algún criterio que usaste para decidir (aunque no lo hayas llamado «cálculo»).}

\begin{infoband}{milcVerde}


\textbf{En tu cuaderno (Actividad 1 · IDENTIFICA).} Título: «Actividad 1 --- Cuando decidí arriesgar». \textbf{Formato:} lista de 3 a 5 viñetas. \textbf{Extensión:} una viñeta por elemento observado (situación, qué ganabas, qué perdías, criterio de decisión). \textbf{Sabes que terminaste cuando} tu lista muestra, en tus propias palabras, que tomaste una decisión con información incompleta pesando riesgo y recompensa.
\end{infoband}

\puente{Acabas de pensar en la lógica del riesgo desde lo que ya sabes. Ahora ponemos esa lógica en números: la tabla de puntuación de Bebras y lo que te dice sobre cuándo responder.}

\milcestacion{milcTurquesa}{Estación 2 de 3 · Entender}

\begin{softbox}{milcTurquesa}{milcGris}{Lo que vas a entender}
La puntuación de Bebras no es un detalle administrativo: es el \textbf{corazón de la estrategia}. Una vez que entiendes los tres niveles y sus penalizaciones, la mejor jugada en cada situación se deduce sola --- como el jugador de canicas que ya no adivina, sino que calcula. Aquí están las cuatro piezas de esa lógica.
\end{softbox}

\milcpaso{1}{milcTurquesa}{\textbf{Nivel A --- responde siempre:} acierto $+6$, error $0$, blanco $0$. No hay penalización. El peor resultado posible al responder es el mismo que dejar en blanco. Adivinar al azar entre 4 opciones da un valor esperado de $+1{,}5$ en promedio. Nunca conviene dejar una A sin respuesta.}
\milcpaso{2}{milcTurquesa}{\textbf{Nivel B --- descarta antes de arriesgar:} acierto $+9$, error $-2$, blanco $0$. Adivinar al azar sin descartar da un valor esperado de $-0{,}5$ (pierdes en promedio). Pero si puedes eliminar 2 de 4 opciones, el valor esperado sube a $+3{,}5$: ahí sí vale arriesgar. Si no descartaste ninguna, deja en blanco.}
\milcpaso{3}{milcTurquesa}{\textbf{Nivel C --- máxima cautela:} acierto $+12$, error $-4$, blanco $0$. Sin descartar, adivinar al azar da $-1$ en promedio (muy malo). Si descartaste 2 de 4, el valor esperado es $+4$ (positivo: arriesga). Descarta primero, decide después.}
\milcpaso{4}{milcTurquesa}{\textbf{Orden de ataque --- asegura primero lo seguro:} empieza por las A (puntos sin riesgo), luego aborda las B con criterio, cierra con las C que hayas podido descartar. Si te trabas en una tarea, márcala y pasa; volver con la mente despejada vale más que tres minutos atascado. El reloj es parte del problema.}

\begin{drawbox}{milcTurquesa}{Cómo decidir en cada nivel}
\textbf{1. Lee el nivel} de la tarea (A, B o C) antes de leerla completa. \\[1mm] \textbf{2. Lee la tarea} y subraya solo los datos que cambian la respuesta (abstracción). \\[1mm] \textbf{3. Descarta} las opciones que puedes eliminar con seguridad. \\[1mm] \textbf{4. Aplica la regla:} A $\to$ responde siempre; B o C con $\leq 1$ opción descartada $\to$ blanco; B o C con $\geq 2$ opciones descartadas $\to$ arriesga. \\[1mm] \textbf{5. Si no avanzas en 3 minutos, pasa} y vuelve al final. El tiempo es un recurso, no un castigo.
\end{drawbox}

\begin{softbox}{milcMostaza}{milcGris}{Errores comunes y su corrección}
\textbf{Responder por impulso en C sin descartar:} adivinar al azar en C da $-1$ en promedio; un punto perdido que nunca se recupera. \textbf{Dejar en blanco las A por «no estar seguro del todo»:} en A el error no resta nada, así que siempre conviene intentarlo. \textbf{Empezar por la más difícil para «ganar más»:} una C no resuelta en 5 minutos te quita tiempo de tres A seguras.
\end{softbox}

\begin{infoband}{milcTurquesa}


\textbf{En tu cuaderno (Actividad 2 · ANALIZA).} Título: «Actividad 2 --- La tabla que decide». \textbf{Formato:} tabla de 4 columnas (nivel · acierto · error · cuándo respondo). \textbf{Extensión:} las tres filas (A, B, C) más una fila de «orden de ataque». \textbf{Sabes que terminaste cuando} un compañero podría usar tu tabla para decidir qué hacer en cualquier situación del concurso sin pedirte explicación.
\end{infoband}

\puente{Ya distingues las reglas y el razonamiento que hay detrás. Es hora de usarlos: vas a enfrentar un escenario real de concurso y decidir, con los números en la mano, cuál es la mejor jugada.}

\milcestacion{milcMagenta}{Estación 3 de 3 · Hacer}

\begin{softbox}{milcMagenta}{milcGris}{Cómo vas a construir tu producto}
La tabla no es adorno: se usa bajo presión. Vas a enfrentar un \textbf{escenario real} de los últimos minutos del concurso, calcular el valor esperado de cada opción y tomar la decisión que maximiza tus puntos. Después escribirás tu plan de ataque personal.
\end{softbox}

\milcpaso{1}{milcMagenta}{\textbf{Lee el escenario completo} antes de decidir: qué tareas hay, qué nivel tiene cada una, cuánto tiempo queda.}
\milcpaso{2}{milcMagenta}{\textbf{Aplica la regla de descarte:} para cada tarea, determina cuántas opciones puedes eliminar. Ese número cambia la decisión.}
\milcpaso{3}{milcMagenta}{\textbf{Calcula el valor esperado} de las opciones que no dominas: multiplica probabilidad por puntuación y suma. El número te dice si vale la pena arriesgar.}
\milcpaso{4}{milcMagenta}{\textbf{Ordena las tareas} de mayor a menor certeza: primero las que ya sabes (A sin duda), luego las que puedes descartar (B o C con $\geq 2$ descartadas), al final las que no puedes descartar (blanco).}

\begin{drawbox}{milcMagenta}{El reto: los últimos 10 minutos}
\checkbox \textbf{Escenario.} Te quedan \textbf{10 minutos}. Tienes 3 tareas sin responder: (i) una \textbf{A} que no has leído, (ii) una \textbf{B} que leíste y no puedes descartar ninguna opción, (iii) una \textbf{C} que leíste y descartaste 2 de 4 opciones. \\[2mm]
\checkbox \textbf{Paso 1.} Calcula el valor esperado de responder al azar la B sin descartar: $(1/4) \times 9 + (3/4) \times (-2)$. \\[2mm]
\checkbox \textbf{Paso 2.} Calcula el valor esperado de arriesgar la C con 2 opciones descartadas: $(1/2) \times 12 + (1/2) \times (-4)$. \\[2mm]
\checkbox \textbf{Paso 3.} Decide el orden en que atacas las 3 tareas y justifica cada decisión. \\[2mm]
\checkbox Escribe tu conclusión en una frase: ¿qué decisión tomas en la B y por qué?
\end{drawbox}

\begin{softbox}{milcMagenta}{milcGris}{Plantilla de trabajo}
\textbf{Solución paso a paso.} \textbf{Paso 1 (B sin descartar):} $(1/4) \times 9 + (3/4) \times (-2) = 2{,}25 - 1{,}5 = -0{,}5$. Valor esperado negativo: dejar en blanco ($0$) es mejor que adivinar ($-0{,}5$). \textbf{Decisión B: blanco.} \textbf{Paso 2 (C con 2 descartadas):} $(1/2) \times 12 + (1/2) \times (-4) = 6 - 2 = +4$. Valor esperado positivo: arriesgar es mejor que dejar en blanco. \textbf{Decisión C: arriesga.} \textbf{Paso 3 (orden de ataque):} primero la A (puntos seguros, sin penalización), luego la C (valor esperado $+4$, positivo), por último la B en blanco (valor esperado $0$, mejor que $-0{,}5$). \textbf{Respuesta correcta: A primero, C con riesgo calculado, B en blanco.}
\end{softbox}

\begin{infoband}{milcMagenta}


\textbf{En tu cuaderno (Actividad 3 · EVALÚA).} Título: «Actividad 3 --- Mi mejor jugada». \textbf{Formato:} tabla de 3 columnas (tarea · decisión · razón con valor esperado) más una frase de conclusión. \textbf{Extensión:} 3 filas más 2 renglones de conclusión. \textbf{Sabes que terminaste cuando} tu justificación usa los números del valor esperado, no solo la intuición.
\end{infoband}

\puente{Tomaste la decisión estratégica en el escenario. Ahora tradúcela a un plan de ataque propio: tu guía personal para el día del concurso, redactada con los criterios que ya conoces.}

\milcestacion{milcMostaza}{Producto de la sesión}
\textbf{Tabla de decisión estratégica del escenario más plan de ataque personal para los 45 minutos}.
\begin{softbox}{milcMostaza}{milcGris}{Criterios de calidad}
(1) Usaste los valores $+6/0$, $+9/-2$, $+12/-4$ de forma correcta en tu razonamiento y los escribiste en la tabla. (2) Calculaste el valor esperado de la B sin descartar ($-0{,}5$) y lo usaste para justificar la decisión de dejar en blanco. (3) Calculaste el valor esperado de la C con 2 opciones descartadas ($+4$) y lo usaste para justificar la decisión de arriesgar. (4) El orden de ataque propuesto es A $\to$ C $\to$ B (blanco) y está justificado con los valores esperados, no con lo que «se siente». (5) Tu plan de ataque personal para los 45 minutos incluye una regla explícita sobre cuándo dejar en blanco (basada en la puntuación) y empieza siempre por las tareas A.
\end{softbox}

\puente{Terminaste el plan. Detente a mirar qué cambió en ti --- no la nota, sino tu manera de decidir cuando la información es incompleta y el tiempo aprieta.}

\milcestacion{milcVino}{Evaluación liberadora · cinco dimensiones}

\begin{softbox}{milcVino}{milcGris}{Sentido de la evaluación}
Evaluar no es solo revisar una nota. Es reconocer qué aprendí, qué cambió en mí, qué emociones manejé, cómo me relaciono con la ciudad y la comunidad, y qué le dejaré a quien viene detrás.
\end{softbox}

\textbf{Mi reflexión liberadora en cinco dimensiones.} Completa cada frase iniciada:

\small
\begin{tabularx}{\textwidth}{>{\bfseries\RaggedRight\arraybackslash}p{3.4cm}Y>{\RaggedRight\arraybackslash}p{4.6cm}}
\rowcolor{milcVino}\color{white}Dimensión & \color{white}\bfseries Mi respuesta (frame starter) & \color{white}\bfseries Algo que descubrí\\
1. Personal & Lo que más cambió en mí fue \linead{6.5cm} & \linead{4.2cm}\\
2. Emocional & La emoción más fuerte fue \linead{2.5cm} y la regulé \linead{3cm} & \linead{4.2cm}\\
3. Ciudadana & En democracia esto importa porque \linead{6cm} & \linead{4.2cm}\\
4. Local (Cartago) & En mi barrio sería útil para \linead{6.5cm} & \linead{4.2cm}\\
5. Intergeneracional & A mi abuela/abuelo le diría \linead{6.5cm} & \linead{4.2cm}\\
\end{tabularx}
\normalsize

\begin{infoband}{milcVino}
\textbf{Para la próxima sustentación.} Vuelve a esta tabla en seis meses. Verás que las respuestas cambian — no porque lo aprendido cambie, sino porque tú cambias. La evaluación liberadora es una conversación contigo a través del tiempo.
\end{infoband}


\puente{Cerramos pensando en grande: tres voces para entender qué significa, para ti y para tu comunidad, aprender a decidir bien bajo presión.}

\milcestacion{milcGradeDark}{Triángulo de pensamiento · cierre filosófico}

\begin{softbox}{milcVino}{milcGris}{Dussel · lente del nosotros}
\textit{La sabiduría popular no es superstición: es razón práctica destilada por siglos de experiencia colectiva.}

Los juegos de canicas y trompo enseñaban estrategia sin libros, solo con práctica y comunidad. La inteligencia táctica que tus mayores te transmitieron en la calle --- leer el terreno, medir el riesgo, decidir con calma --- es exactamente la misma que Bebras pone a prueba con números. El saber popular ya era pensamiento computacional.

\textbf{¿Qué otros saberes de tu barrio o vereda son, en realidad, inteligencia práctica que la escuela aún no ha reconocido como tal?}
\end{softbox}
\linead{\linewidth}\\[1mm]
\linead{\linewidth}

\begin{softbox}{milcMostaza}{milcGris}{Estoicismo · Marco Aurelio · cuidado interior}
\textit{Si no está en tu poder cambiar el resultado, al menos está en tu poder elegir cómo afrontar la situación.}

En el concurso no puedes controlar qué tareas aparecen ni cuánto tiempo hay. Sí puedes controlar cómo decides bajo presión: con calma y estrategia, sin dejarte llevar por el afán de responder todo. La ecuanimidad no es rendirse: es decidir con la cabeza fría.

\textbf{¿En qué otras situaciones de tu vida --- un examen, una discusión, una decisión difícil --- aplicarías la misma lógica: controlar lo que sí está en tu mano y soltar lo que no?}
\end{softbox}
\linead{\linewidth}\\[1mm]
\linead{\linewidth}

\begin{softbox}{milcTurquesa}{milcGris}{Floridi · lente de la infoesfera}
\textit{Tomar decisiones racionales bajo incertidumbre es una de las competencias más valiosas en la era de la información.}

Los algoritmos de inteligencia artificial también «apuestan» con probabilidades cuando no tienen información completa. Tu decisión de responder o dejar en blanco una tarea C es, en esencia, la misma operación que hace una máquina al predecir con datos incompletos: calcular el valor esperado y elegir la opción que lo maximiza.

\textbf{¿Cómo se parece tu decisión estratégica en Bebras a la que toma una máquina al clasificar o predecir algo con información parcial?}
\end{softbox}
\linead{\linewidth}\\[1mm]
\linead{\linewidth}

\begin{infoband}{milcNegro}
\textbf{Nota para el docente.} Las tres formulaciones del triángulo son la síntesis del autor de la guía, no citas textuales de los pensadores. Si vas a citarlos en clase, remite a la obra original.
\end{infoband}

\milcestacion{milcVino}{Mi autoevaluación · marca una casilla por fila}

{\small
\setlength{\extrarowheight}{2.2pt}
\begin{tabularx}{\textwidth}{>{\RaggedRight\arraybackslash\bfseries}p{3.0cm}YYY}
\rowcolor{milcVino}\color{white}\bfseries Criterio & \color{white}\bfseries Lo logré & \color{white}\bfseries En proceso & \color{white}\bfseries Necesito apoyo\\[.6mm]
Comprensión del tema & \milcopcion{Explico las ideas con mis palabras.} & \milcopcion{Reconozco las ideas, falta precisión.} & \milcopcion{Necesito repasar lo básico.}\\[.8mm]
\rowcolor{milcGris}Pensamiento computacional & \milcopcion{Usé los pilares al armar mi producto.} & \milcopcion{Usé algunos pilares.} & \milcopcion{Necesito releer la estación 2.}\\[.8mm]
Aplicación y producto & \milcopcion{Mi producto cumple los criterios.} & \milcopcion{Está armado, con ajustes pendientes.} & \milcopcion{Aún está en construcción.}\\[.8mm]
\rowcolor{milcGris}Reflexión 5 dimensiones & \milcopcion{Respondí cada una con honestidad.} & \milcopcion{Respondí algunas con detalle.} & \milcopcion{Mis respuestas son muy generales.}\\[.8mm]
Triángulo de pensamiento & \milcopcion{Conecté las 3 voces con mi trabajo.} & \milcopcion{Conecté al menos una.} & \milcopcion{No las relaciono con mi proceso.}\\[.8mm]
\end{tabularx}}

\vspace{3mm}
\textbf{Hoy entendí que:}\\[1mm]
\linead{\linewidth}\\[1mm]
\linead{\linewidth}

\textbf{Mi compromiso para la próxima sesión es:}\\[1mm]
\linead{\linewidth}\\[1mm]
\linead{\linewidth}

\begin{softbox}{milcMostaza}{milcGris}{Cita de despedida · Marco Aurelio}
\textit{``El obstáculo en el camino se vuelve el camino. Lo que estorba al camino se vuelve camino.''}\\[1mm]
Cada error de hoy, bien observado, se convierte en pieza del camino que pisarás mañana.
\end{softbox}

\small
\begin{tabularx}{\textwidth}{>{\bfseries}p{3.6cm}Y}
\rowcolor{milcGradeDark!12}Nombre completo: & \linead{10cm}\\
Fecha: & \linead{4cm} \quad Curso/grupo: \linead{4cm}\\
Mi firma: & \linead{9cm}\\
\end{tabularx}
\normalsize

\begin{infoband}{milcGradeDark}
\textbf{Para la próxima sesión.} Pon a prueba tu plan de ataque en el simulacro cronometrado del Momento~9: 45 minutos, tareas reales, estrategia aplicada. Antes de llegar, repasa tu tabla de decisión y la regla que escribiste sobre cuándo dejar en blanco. Si tienes un familiar o amigo que juegue dominó, ajedrez o cualquier juego de mesa, pregúntale cómo decide cuándo arriesgar: comparte lo que encontraste con el grupo.
\end{infoband}

\end{document}
